% GENERATED FILE. Edit canonical lessons, then run npm run generate:books.
% canonical-source-hash: fnv1a64:781456f5b2f46215
% canonical-lessons: SA-C341-prayojanam, SA-C341-uddesyam, SA-C341-upayogah, SA-C341-prerayati, SA-C341-karayati

\chapter{What It Is For}
\label{ch:sa-what-it-is-for}
\hlchaptermodality{341}

\begin{chapteropening}
\textbf{By the end of this chapter:} \emph{I can say what something is for and why people do it.}

Complete the last lesson of chapter 341: I can say what something is for and why people do it.
\end{chapteropening}

\section[\texorpdfstring{\sk{प्रयोजनम्} (prayojanam)}{प्रयोजनम् (prayojanam)}]{\sk{प्रयोजनम्} (prayojanam) — a purpose}

\label{lesson:SA-C341-prayojanam}



\begin{quote}
Before the new one: say the Sanskrit for comes about, then the Sanskrit for gives rise to.
\end{quote}

\subsection*{You'll want to know: \sk{प्रयोजनम्}}
\textbf{\sk{प्रयोजनम्}} — \emph{prayojanam} — \textquotedblleft{}a purpose\textquotedblright{}.

\begin{quote}
\sk{एतस्य} \sk{कार्यस्य} \sk{प्रयोजनं} \sk{किम्}?
\end{quote}

\begin{quote}
\emph{etasya kāryasya prayojanaṃ kim?}
\end{quote}

\begin{quote}
What is the purpose of this work?
\end{quote}

A neuter noun.

\begin{cousinweb}
\textbf{\sk{प्रयोजनम्}} (\emph{prayojanam}) is built on \emph{pra} and the root √yuj, \textquotedblleft{}to join, to use\textquotedblright{}; MW gives \emph{pra-yojana} the senses \textquotedblleft{}occasion, motive, purpose, aim\textquotedblright{} (MW).
\end{cousinweb}

\begin{grammarlens}[title={What you've built}]
One more word for giving reasons and opinions.
\end{grammarlens}

\subsection*{Your turn}
\begin{itemize}
\raggedright
  \item \emph{Say it:} \emph{prayojanam}
  \item \emph{Say it:} \emph{prayojanam}, once more
  \item \emph{Say it:} \emph{etasya kāryasya prayojanaṃ kim?} — What is the purpose of this work?
  \item \emph{Recall:} say the Sanskrit for comes about, then the Sanskrit for gives rise to, then say \emph{prayojanam} again
\end{itemize}

\begin{tcolorbox}[breakable,colback=teal!4,colframe=teal!35!black,title={Before you move on}]
What does \sk{प्रयोजनम्} mean? (\textquotedblleft{}A purpose\textquotedblright{}.) Say it once more.
\end{tcolorbox}
\section[\texorpdfstring{\sk{उद्देश्यम्} (uddeśyam)}{उद्देश्यम् (uddeśyam)}]{\sk{उद्देश्यम्} (uddeśyam) — an aim}

\label{lesson:SA-C341-uddesyam}



\begin{quote}
Before the new one: say the Sanskrit for gives rise to, then the Sanskrit for a purpose.
\end{quote}

\subsection*{You'll want to know: \sk{उद्देश्यम्}}
\textbf{\sk{उद्देश्यम्}} — \emph{uddeśyam} — \textquotedblleft{}an aim\textquotedblright{}.

\begin{quote}
\sk{मम} \sk{उद्देश्यं} \sk{स्पष्टम्} \sk{अस्ति।}
\end{quote}

\begin{quote}
\emph{mama uddeśyaṃ spaṣṭam asti.}
\end{quote}

\begin{quote}
My aim is clear.
\end{quote}

A neuter noun.

\begin{cousinweb}
\textbf{\sk{उद्देश्यम्}} (\emph{uddeśyam}) is built on \emph{ud}, \textquotedblleft{}up\textquotedblright{}, and the root √diś, \textquotedblleft{}to point out\textquotedblright{}; MW gives \emph{ud-deśya} as \textquotedblleft{}anything which one has in view\textquotedblright{}, and \textquotedblleft{}an aim\textquotedblright{} is the sense modern Sanskrit gives it (MW).
\end{cousinweb}

\begin{grammarlens}[title={What you've built}]
One more word for giving reasons and opinions.
\end{grammarlens}

\subsection*{Your turn}
\begin{itemize}
\raggedright
  \item \emph{Say it:} \emph{uddeśyam}
  \item \emph{Say it:} \emph{uddeśyam}, once more
  \item \emph{Say it:} \emph{mama uddeśyaṃ spaṣṭam asti.} — My aim is clear.
  \item \emph{Recall:} say the Sanskrit for gives rise to, then the Sanskrit for a purpose, then say \emph{uddeśyam} again
\end{itemize}

\begin{tcolorbox}[breakable,colback=teal!4,colframe=teal!35!black,title={Before you move on}]
What does \sk{उद्देश्यम्} mean? (\textquotedblleft{}An aim\textquotedblright{}.) Say it once more.
\end{tcolorbox}
\section[\texorpdfstring{\sk{उपयोगः} (upayogaḥ)}{उपयोगः (upayogaḥ)}]{\sk{उपयोगः} (upayogaḥ) — a use}

\label{lesson:SA-C341-upayogah}



\begin{quote}
Before the new one: say the Sanskrit for a purpose, then the Sanskrit for an aim.
\end{quote}

\subsection*{You'll want to know: \sk{उपयोगः}}
\textbf{\sk{उपयोगः}} — \emph{upayogaḥ} — \textquotedblleft{}a use\textquotedblright{}.

\begin{quote}
\sk{एतस्य} \sk{यन्त्रस्य} \sk{उपयोगः} \sk{कः}?
\end{quote}

\begin{quote}
\emph{etasya yantrasya upayogaḥ kaḥ?}
\end{quote}

\begin{quote}
What is the use of this machine?
\end{quote}

A masculine noun.

\begin{cousinweb}
\textbf{\sk{उपयोगः}} (\emph{upayogaḥ}) is built on \emph{upa}, \textquotedblleft{}towards\textquotedblright{}, and the root √yuj, \textquotedblleft{}to join\textquotedblright{}; MW gives \emph{upa-yoga} the senses \textquotedblleft{}employment, use, application\textquotedblright{} (MW).
\end{cousinweb}

\begin{grammarlens}[title={What you've built}]
One more word for giving reasons and opinions.
\end{grammarlens}

\subsection*{Your turn}
\begin{itemize}
\raggedright
  \item \emph{Say it:} \emph{upayogaḥ}
  \item \emph{Say it:} \emph{upayogaḥ}, once more
  \item \emph{Say it:} \emph{etasya yantrasya upayogaḥ kaḥ?} — What is the use of this machine?
  \item \emph{Recall:} say the Sanskrit for a purpose, then the Sanskrit for an aim, then say \emph{upayogaḥ} again
\end{itemize}

\begin{tcolorbox}[breakable,colback=teal!4,colframe=teal!35!black,title={Before you move on}]
What does \sk{उपयोगः} mean? (\textquotedblleft{}A use\textquotedblright{}.) Say it once more.
\end{tcolorbox}
\section[\texorpdfstring{\sk{प्रेरयति} (prerayati)}{प्रेरयति (prerayati)}]{\sk{प्रेरयति} (prerayati) — urges on; inspires}

\label{lesson:SA-C341-prerayati}



\begin{quote}
Before the new one: say the Sanskrit for an aim, then the Sanskrit for a use.
\end{quote}

\subsection*{You'll want to know: \sk{प्रेरयति}}
\textbf{\sk{प्रेरयति}} — \emph{prerayati} — \textquotedblleft{}urges on; inspires\textquotedblright{}.

\begin{quote}
\sk{आचार्यः} \sk{शिष्यान्} \sk{प्रेरयति।}
\end{quote}

\begin{quote}
\emph{ācāryaḥ śiṣyān prerayati.}
\end{quote}

\begin{quote}
The preceptor inspires his pupils.
\end{quote}

Past (he …): \sk{सः} \sk{मां} \sk{प्रेरितवान्}, \emph{saḥ māṃ preritavān}.

\begin{cousinweb}
\textbf{\sk{प्रेरयति}} (\emph{prerayati}) is built on \emph{pra}, \textquotedblleft{}forth\textquotedblright{}, and the root √īr, \textquotedblleft{}to set in motion\textquotedblright{}, in its causative form; MW gives \emph{prerita} \textquotedblleft{}urged, impelled, sent\textquotedblright{} (MW).
\end{cousinweb}

\begin{grammarlens}[title={What you've built}]
One more word for giving reasons and opinions.
\end{grammarlens}

\subsection*{Your turn}
\begin{itemize}
\raggedright
  \item \emph{Say it:} \emph{prerayati}
  \item \emph{Say it:} \emph{prerayati}, once more
  \item \emph{Say it:} \emph{ācāryaḥ śiṣyān prerayati.} — The preceptor inspires his pupils.
  \item \emph{Recall:} say the Sanskrit for an aim, then the Sanskrit for a use, then say \emph{prerayati} again
\end{itemize}

\begin{tcolorbox}[breakable,colback=teal!4,colframe=teal!35!black,title={Before you move on}]
What does \sk{प्रेरयति} mean? (\textquotedblleft{}Urges on; inspires\textquotedblright{}.) Say it once more.
\end{tcolorbox}
\section[\texorpdfstring{\sk{कारयति} (kārayati)}{कारयति (kārayati)}]{\sk{कारयति} (kārayati) — has something done}

\label{lesson:SA-C341-karayati}



\begin{quote}
Before the new one: say the Sanskrit for a use, then the Sanskrit for urges on.
\end{quote}

\subsection*{You'll want to know: \sk{कारयति}}
\textbf{\sk{कारयति}} — \emph{kārayati} — \textquotedblleft{}has something done\textquotedblright{}.

\begin{quote}
\sk{राजा} \sk{मन्दिरं} \sk{कारयति।}
\end{quote}

\begin{quote}
\emph{rājā mandiraṃ kārayati.}
\end{quote}

\begin{quote}
The king has a temple built.
\end{quote}

Past (he …): \sk{सः} \sk{गृहं} \sk{कारितवान्}, \emph{saḥ gṛhaṃ kāritavān}.

\begin{cousinweb}
\textbf{\sk{कारयति}} (\emph{kārayati}) comes from the root √kṛ, \textquotedblleft{}to do, to make\textquotedblright{}, in its causative form; MW gives \emph{kārita} \textquotedblleft{}caused to be made or done\textquotedblright{} (MW).
\end{cousinweb}

\begin{grammarlens}[title={What you've built}]
One more word for giving reasons and opinions.
\end{grammarlens}

\subsection*{Your turn}
\begin{itemize}
\raggedright
  \item \emph{Say it:} \emph{kārayati}
  \item \emph{Say it:} \emph{kārayati}, once more
  \item \emph{Say it:} \emph{rājā mandiraṃ kārayati.} — The king has a temple built.
  \item \emph{Recall:} say the Sanskrit for a use, then the Sanskrit for urges on, then say \emph{kārayati} again
\end{itemize}

\begin{tcolorbox}[breakable,colback=teal!4,colframe=teal!35!black,title={Before you move on}]
What does \sk{कारयति} mean? (\textquotedblleft{}Has something done\textquotedblright{}.) Say it once more.
\end{tcolorbox}
